\section{Contexto del taller y protocolo de consentimiento informado}

\subsection{Prop\'osito y alcance del ensayo}
Este documento recoge el dise\~no preliminar y el ensayo controlado de la entrevista estructurada y de la sesi\'on de \emph{design thinking} para el Proyecto Fin de Curso (PFC), orientados a la operaci\'on real de establecimientos de suministros el\'ectricos en el cant\'on Quevedo. La combinaci\'on de entrevista estructurada y herramientas de empat\'ia permite comparar respuestas de forma consistente y detectar necesidades operativas que los usuarios no verbalizan de forma directa \cite{pohl2025,ireb2026}.

El ensayo se ejecuta bajo modalidad de juego de roles dentro del aula, tomando como parte interesada prioritaria al \textbf{Cajero / Vendedor de Mostrador (SH-01)}, dado que este rol concentra la atenci\'on directa en horas pico, la consulta de inventario y la asesor\'ia t\'ecnica sobre compatibilidad de materiales el\'ectricos.

\subsection{Modelo de consentimiento informado (LOPDP Ecuador)}
En cumplimiento del C\'odigo de \'Etica de la ACM \cite{acm2018} y de la Ley Org\'anica de Protecci\'on de Datos Personales (LOPDP) del Ecuador \cite{lopdp2021}, toda sesi\'on de elicitaci\'on requiere consentimiento libre, espec\'ifico, informado e inequ\'ivoco del participante. Ning\'un dato personal de terceros (n\'umeros de c\'edula, tel\'efonos, correos personales, grabaciones de audio o actas firmadas) se almacena en el repositorio p\'ublico de GitHub; dichas evidencias se resguardan en la carpeta institucional privada (\texttt{@uteq.edu.ec}) compartida \'unicamente con el docente, y en el repositorio se publicar\'a el manifiesto con la huella criptogr\'afica SHA-256 correspondiente.

\begin{tcolorbox}[colback=boxbg,colframe=midblue,boxrule=0.8pt,arc=2mm,breakable,title={\textbf{Plantilla de Consentimiento Informado para Registro de Audio (ACTA-CONS-01)}}]
\textbf{Proyecto:} Plataforma de Ventas con M\'odulo de Recomendaci\'on basado en IA para Establecimientos de Suministros El\'ectricos en Quevedo.\\
\textbf{Asignatura:} Ingenier\'ia de Requerimientos (ISR-401) --- Universidad T\'ecnica Estatal de Quevedo.\\[4pt]
\textbf{1. Finalidad del registro:} Comprender el flujo diario de trabajo en mostrador, los tiempos de atenci\'on y los criterios de asesor\'ia t\'ecnica con fines estrictamente acad\'emicos, sin evaluar el desempe\~no laboral individual.\\[3pt]
\textbf{2. Voluntariedad y revocaci\'on:} La participaci\'on es libre y voluntaria. El participante puede omitir cualquier pregunta, solicitar una pausa o retirar su consentimiento en cualquier momento sin consecuencia alguna.\\[3pt]
\textbf{3. Anonimizaci\'on y custodia:} Los nombres reales del participante y del local comercial se reemplazan por c\'odigos de rol (por ejemplo, \texttt{SH-01}). El archivo de audio se conserva exclusivamente en almacenamiento institucional privado UTEQ bajo los principios de confidencialidad y minimizaci\'on de datos de la LOPDP \cite{lopdp2021}.\\[6pt]
\noindent\begin{tabularx}{\linewidth}{|Y|Y|}
\hline
\textbf{C\'odigo de sesi\'on:} \texttt{GA-TL-02-ENS-01} & \textbf{Fecha:} \_\_\_ / \_\_\_ / 2026\\
\hline
\textbf{Rol representado:} Cajero / Vendedor (SH-01) & \textbf{Autoriza grabaci\'on de audio:} [ \ ] S\'i \quad [ \ ] No\\
\hline
\textbf{Nombre del entrevistador:} \rule{0pt}{14pt} & \textbf{Firma / confirmaci\'on verbal:} \rule{0pt}{14pt}\\
\hline
\end{tabularx}
\end{tcolorbox}
